\documentclass[9pt,a4paper]{extarticle}
\usepackage{parskip}

\usepackage[utf8]{inputenc}
\usepackage[T1]{fontenc}
\usepackage[brazil]{babel}
\usepackage{amsmath, amssymb, amsthm}
\usepackage{geometry}
\usepackage{xcolor}
\usepackage{hyperref}
\usepackage{booktabs}
\usepackage{enumitem}
\usepackage{microtype}
\usepackage{csquotes}
\usepackage{mdframed}
\usepackage{pgffor}
\usepackage{tabto}

\newcommand{\answerlines}[1]{%
  \par\vspace{4pt}%
  \foreach \n in {1,...,#1}{%
    \noindent\makebox[\linewidth]{\dotfill}\\[6pt]%
  }%
}

\definecolor{important}{RGB}{255, 100, 100}
\definecolor{notice}{RGB}{0, 102, 204}
\definecolor{headerblue}{RGB}{152, 38, 73}
\definecolor{accentblue}{RGB}{113, 162, 182}

\newcommand{\impt}[1]{\textcolor{important}{\textbf{#1}}}
\newcommand{\ntc}[1]{\textcolor{notice}{#1}}

\setlist{nosep}
\geometry{margin=2.5cm, top=1.5cm}

% Caixa para observações/gabarito
\newmdenv[
  linecolor=headerblue,
  backgroundcolor=headerblue!6,
  linewidth=1.2pt,
  innerleftmargin=10pt,
  innerrightmargin=10pt,
  innertopmargin=8pt,
  innerbottommargin=8pt,
  skipabove=6pt,
  skipbelow=6pt
]{notebox}

\title{%
  \vspace{-1cm}%
  \textcolor{headerblue}{\textbf{Prova 1}}\\[4pt]
  \large Apresentação dos Cursos de CC e IA --- DECOM/UFOP
}
\date{}

\begin{document}
\maketitle
\vspace{-1cm}

\begin{notebox}
  \textbf{Instruções.} Responda à questão a seguir de forma clara, objetiva e bem fundamentada.
  Não basta descrever o que está na grade curricular, espera-se \impt{análise crítica}.
\end{notebox}

\vspace{0.5cm}

% ------------------------------------------------------------------
\begin{enumerate}

  \item Os cursos de Ciência da Computação e Inteligência Artificial compartilham
  disciplinas como Estruturas de Dados, POO, Teoria dos Grafos e Algoritmos, além
  do mesmo corpo docente, laboratórios e acesso ao PPGCC.
  Um estudante argumenta que, diante disso, manter dois cursos separados é
  \ntc{redundante e ineficiente}. Avalie criticamente esse argumento, considerando
  tanto os pontos que o \impt{sustentam} quanto os que o \impt{contradizem}.

\end{enumerate}

\answerlines{75}

\end{document}